\begin{abstract}
We suggest a new proof of quantum fault tolerance based on the alternative proofs of Toom's theorem \cite{Berman} (see also \cite{gacs2021newversiontoomsproof}). The fault tolerance construction is relative to the circuit model, subjected to local stochastic noise, where at each time unit only a constant depth of non-noisy classical computation is allowed. Our main error correction code is the $(d,d^{\prime})$-Toric code, namely qubits are set on the $d^{\prime}$-dimensional faces of the $d$-dimensional Toric. As byproducts of the study, we provide and prove the following:
  \begin{enumerate}
    \item  An explicit description of the sweep-rule —a generalization of Toom's rule for topological codes—for any dimension greater than two. In doing so, it answers affirmatively to the conjecture made by Kubica and Preskill \cite{Kubica_2019}.
    \item A fault tolerant construction in the circuit model that both uses the $(d,d^{\prime})$-Toric code and has a constant-depth correction cycle, where $d^{\prime}, d - d^{\prime} > 2$. Until now, the use of a topological code as the primary encoding and a constant-depth correction cycle were only achieved separately; this is the first work to achieve them simultaneously.
    \item A novel fault tolerance that doesn't introduce an asymptotic depth overhead per Clifford gate. In particular, there is a threshold $p_{0} \in (0,1)$ such that for a noise rate $p < p_{0}$, and any $n$ large enough, the construction takes a circuit with total depth $\mu$ and a $\mathbf{T}$-depth $\mu_{\mathbf{T}}$, and results a circuit that simulates the given with probability: 
      \begin{equation*}
        \begin{split}
          \Theta ( \text{Poly(circuit's size)} ) p^{\Theta(n)}   
        \end{split}
      \end{equation*}
      and has a depth $\Theta(\mu) + \Theta( \mu_{\mathbf{T}} n^{d+1} \log n ) $. 
  \end{enumerate}
Specifically, all these results hold for the 4D-Toric code.
\end{abstract}
